\documentclass[12pt]{article}
\usepackage[utf8]{inputenc}
\usepackage[T1]{fontenc}
\usepackage[margin=1in]{geometry}
\usepackage{setspace}
\usepackage{xcolor}
\usepackage[colorlinks=true, urlcolor=blue!60!black, linkcolor=blue!60!black]{hyperref}
\setlength{\parindent}{0em}
\setlength{\parskip}{0.45em}
\renewcommand{\baselinestretch}{1.0}
\begin{document}
\begin{center}
{\Large \textit{Stuart v. Laird}, 5 U.S. (1 Cranch) 299 (1803)\par}
\vspace{0.3em}
{\normalsize Case Brief --- One-Page Edition\par}
\end{center}

\textbf{Memory Jogger.} Congress may transfer cases between reorganized lower courts, and practice settled circuit riding (\href{https://supreme.justia.com/cases/federal/us/5/299/}{Justia}).

\textbf{Facts.} Laird's debt judgment against Stuart, won in a court created by the Federalists' 1801 ``Midnight Judges Act,'' was transferred after the Republican Congress repealed the Act in 1802 and restored circuit riding. Marshall, on circuit, entered judgment on Stuart's bond.

\textbf{Procedural History.} Stuart brought a writ of error challenging the transfer and circuit riding.

\textbf{Votes.} Affirmed 5--0: Paterson for the Court; Marshall did not participate.

\textbf{Holding.} Congress may transfer pending cases among inferior courts; long practice validates Justices sitting on circuit without new commissions.

\textbf{Reasoning.} The power to create inferior tribunals includes the power to move causes between them. Circuit riding since 1789 was a ``contemporary interpretation of the most forcible nature.'' The Court never addressed the ousted judges' Good Behavior claim; Marshall and Chase privately doubted circuit riding, but the Justices closed ranks amid Jeffersonian threats of impeachment.

\textbf{Examples.} \emph{Same-Side:} splitting the Ninth Circuit. \emph{Opposite-Side:} abolishing one judge's seat to oust her. \emph{Fence-Sitter:} eighteen-year limits for Justices.

\textbf{Critique.} Liquidation theorists such as William Baude see proof that early practice can settle text, but thirteen years of protested practice overrode textual objections, and the tenure question was dodged.

\textbf{Key Quotation.} ``To this objection, which is of recent date, it is sufficient to observe, that practice and acquiescence under it for a period of several years, commencing with the organization of the judicial system, affords an irresistible answer, and has indeed fixed the construction.'' --- Justice Paterson, 5 U.S. (1 Cranch) at 309.

\end{document}
